% ============================================================
% 摘要与关键词
% ============================================================

\begin{abstract}
\noindent 针对问题一，，首先根据舞龙队沿等距螺线盘入的运动特征，建立阿基米德螺线运动模型，并利用弧长关系确定任意时刻龙头前把手的极角。对于非线性方程难以直接求得解析解的问题，采用二分法进行数值求解。在此基础上，以相邻把手间的固定距离为几何约束，利用余弦定理建立相邻把手极角差的非线性方程，并通过二分法求取满足盘入方向的最小正根，逐节递推得到其余把手的极角及平面坐标，从而实现任意时刻舞龙队各把手位置的确定。此后，再通过求导的方式实现各点速度的求解。
\end{abstract}

\vspace{0.5cm}
\noindent \textbf{关键词：} 关键词一，关键词二，关键词三 

